\section{Background}\label{sec:background}

The Assertion Ladder draws on five bodies of work: the emergence of coding agents, architectural
structure and its testing, behaviour and requirement specification, the psychology of learning,
and lean production. This section introduces only what the later sections use.

\subsection{Coding agents and the memoryless contributor}

Code-generating models moved from completing a line \citep{chen2021codex} to resolving whole
issues in real repositories \citep{jimenez2024swebench}. What made the second step possible was
less the model than its harness: an agent that can list files, open them, edit them and run the
tests behaves very differently from one that emits a single completion
\citep{yang2024sweagent}. Tool vendors now converge on a convention for telling such an agent
how a repository works: a Markdown file at the root, loaded at the start of every session
\citep{agentsmd}. That file is the agent's whole inheritance from previous sessions.

Three properties of this contributor matter for what follows. It is \emph{capable}: given a
clear pattern it will reproduce it across layers and languages. It is \emph{memoryless}: nothing
learned in one session survives to the next unless the repository records it. And it is
\emph{fast}: its output can exceed what a person can review line by line, which moves the
person's role from author to reviewer \citep{barke2023grounded} and makes the quality of what
passes review depend on what the reviewer can see.

\subsection{Structure and its enforcement}

Layering is one of the oldest structuring ideas in software \citep{dijkstra1968the}; its
rationale is information hiding, so that a decision likely to change is confined to one module
\citep{parnas1972criteria}. The layered architecture of enterprise applications, with a
presentation layer, a service layer, a data-access layer and a domain model, is catalogued as a
pattern \citep{buschmann1996posa,fowler2002peaa}, and domain-driven design places the domain
model at the centre of attention \citep{evans2003ddd}. Hexagonal and Clean Architecture invert
the dependencies so that the domain depends on nothing and outer adapters depend on ports it
owns \citep{cockburn2005hexagonal,martin2017clean}.

Structure described only in prose decays. Architecture fitness functions make an
architectural characteristic executable, so that a build can fail when it is violated
\citep{ford2017evolutionary}; ArchUnit does this for Java dependencies by analysing bytecode in
an ordinary unit test \citep{archunit}. Such rules are the structural core of our third rung.

\subsection{Behaviour and requirements}

Behaviour-driven development writes tests as examples of behaviour in the vocabulary of the
domain, conventionally in the form \emph{given} a context, \emph{when} an event occurs,
\emph{then} an outcome is expected \citep{north2006bdd}; specification by example treats such
examples as the living specification of the system \citep{adzic2011sbe}.

The Easy Approach to Requirements Syntax (EARS) constrains natural-language requirements to a
small set of templates, each distinguished by a keyword \citep{mavin2009ears}:
\emph{ubiquitous} requirements (``The system shall\ldots''), \emph{state-driven}
(``While\ldots''), \emph{event-driven} (``When\ldots''), \emph{optional feature}
(``Where\ldots''), \emph{unwanted behaviour} (``If\ldots, then\ldots'') and \emph{complex}
combinations of these. EARS was designed to remove ambiguity for human readers. Section~\ref{sec:ladder}
argues that its keywords do more work in a repository shared with agents: they classify a rule
in a way that predicts where its code belongs.

\subsection{Learning: worked examples, scaffolding and cognitive load}

Cognitive load theory distinguishes the load a task imposes because of its intrinsic difficulty
from load imposed by how the task is presented, and holds that working memory is small
\citep{sweller1988cognitive}. One of its best-established findings is the worked-example
effect: novices learn a procedure faster by studying a solved example than by solving the
problem themselves \citep{sweller1985worked}. Scaffolding describes support that lets a learner
do what they could not do alone, and that is withdrawn as competence grows \citep{wood1976scaffolding}.

A coding agent is not a novice in the usual sense; it is highly competent in general and
entirely ignorant of this repository in particular. The analogy we use is limited to that
second property. The repository must teach the particular, every session, within a context
window that is large but not unlimited.

\subsection{Lean production and lean software}

Lean production, derived from the Toyota Production System, treats anything that does not add
value for the customer as waste, and builds quality into the process rather than inspecting it
in at the end \citep{ohno1988tps,womack1996lean}. Four of its mechanisms recur in this paper.
\emph{Jidoka} stops the line at the first defect so that it cannot propagate.
\emph{Poka-yoke} designs a fixture or step so that a mistake cannot be made, or is detected at
the moment it is made \citep{shingo1986zqc}. \emph{Standard work} records the current best method
so that it can be followed and improved. \emph{Andon} gives every worker a way to stop and call
for help. Lean software development maps these onto programming practices such as automated
tests, continuous integration and small batches \citep{poppendieck2003lean}, and the
build--measure--learn loop brings the scientific method to product decisions
\citep{ries2011lean}.

\subsection{Patterns}

A pattern names a recurring problem in a context, the forces that make it hard, and a solution
that balances them \citep{alexander1977pattern}. The form was brought to software design by the
Gang of Four \citep{gamma1994design} and to architecture by the POSA series
\citep{buschmann1996posa}. We use it for the Assertion Ladder because the ladder is not an
algorithm or a tool but a way of resolving forces that every team sharing a repository with
agents will meet.
